\lecture{24}
\title{Linear Regression II}

\begin{document}

\subsection{Variance of the LS Solution}

Here's an alternative way of presenting the assumptions
we made about the relationship between feature vectors $X$
and response variables $Y$ from
\href{/18.650/lectures/23/#linear-regression-mle}{Lecture 23}.

\begin{quote}
    \textbf{Assumption.}
    For some true unknown $\beta^*$,
    we have $Y_i = X_i^{\top}\beta^* + \epsilon_i$,
    where $\epsilon_i \sim \mathcal{N}(0, \sigma^2)$
    for some (unknown) constant $\sigma$.
\end{quote}

Under this model, the maximum-likelihood estimator
is $\hat{\beta}^{\mathrm{MLE}} = \hat{\beta}^{\mathrm{LS}}
= (\mathbb{X}^{\top}\mathbb{X})^{-1}\mathbb{X}^{\top}Y$.
But what's the variance of this estimator?
How far does it tend to be off from the true value $\beta^*$?

\begin{remark}
    Recall that the estimator $\hat{\beta}^{\mathrm{LS}}$
    is still an MLE, so all the properties of MLE---consistency,
    asymptotic normality---still hold.
    (But finding the asymptotic variance
    isn't as straightforward as computing the Fisher information.)
\end{remark}

\begin{theorem}{LS Variance}
    We have $\hat{\beta}^{\mathrm{LS}}
    \sim \mathcal{N}_k(\beta^*, \sigma^2(\mathbb{X}^{\top}\mathbb{X})^{-1})$.
\end{theorem}

\begin{proof}
    More generally, it is true that the vector $Y \in \mathbb{R}^n$
    satisfies $Y = \mathbb{X}\beta^* + \epsilon$,
    where $\epsilon \sim \mathcal{N}_n(0, \sigma^2 I_n)$.  Thus,
    \[ \hat{\beta}^{\mathrm{LS}}
    = (\mathbb{X}^{\top}\mathbb{X})^{-1}\mathbb{X}^{\top}(
        \mathbb{X}\beta^* + \epsilon
    )
    = \beta^* + (\mathbb{X}^{\top}\mathbb{X})^{-1}\mathbb{X}^{\top}\epsilon. \]
    The rest of the proof is just multivariate Gaussian properties; recall
    \href{/18.650/lectures/6/#multivariate-gaussian-properties-and-theorems}{Lecture 6},
    or more fundamentally, the properties of
    covariance matrices from \href{/18.650/lectures/5/}{Lecture 5}.

    The mean of the RHS is clearly $\beta^*$,
    whereas the covariance is:
    \begin{align*}
    \Cov\left(
        \left[(\mathbb{X}^{\top}\mathbb{X})^{-1}\mathbb{X}^{\top}\right]
    \epsilon \right)
    & = \left[(\mathbb{X}^{\top}\mathbb{X})^{-1}\mathbb{X}^{\top}\right]
    \Cov(\epsilon)
    \left[(\mathbb{X}^{\top}\mathbb{X})^{-1}\mathbb{X}^{\top}\right]^{\top} \\
    & = \sigma^2 \cdot
    (\mathbb{X}^{\top}\mathbb{X})^{-1}\mathbb{X}^{\top}
    \mathbb{X}(\mathbb{X}^{\top}\mathbb{X})^{-1} \\
    & = \sigma^2 (\mathbb{X}^{\top}\mathbb{X})^{-1}.
    \end{align*}
    Yeah, the computations look terrible,
    but at least it matches what the theorem promised. ~ $\blacksquare$
\end{proof}

\subsection{Linear Regression Confidence Intervals}

We know the variance now,
so surely computing confidence intervals and performing tests
should be easy, right?

There's still a problem: we don't know $\sigma$,
so the variance being $\sigma^2(\mathbb{X}^{\top}\mathbb{X})^{-1}$
still doesn't tell us anything yet.
But---in a similar fashion as bootstrapping---there's a trick
to use the training data to estimate $\sigma$.

\begin{theorem}{Estimating Regression Variance}
    Given the training data and $\hat{\beta}^{\mathrm{LS}}$,
    we may compute the vector
    $\hat{\epsilon} := Y - \mathbb{X}\hat{\beta}^{\mathrm{LS}}$.
    Then the estimator $\hat{\sigma}^2 := \frac{1}{n - k} \|\hat{\epsilon}\|^2$
    is an unbiased estimator for $\sigma^2$.
\end{theorem}

\begin{remark}
    This is a very technically dense proof,
    but the fundamental ideas are relatively straightforward.
    \begin{center}
        \frame{\includegraphics[width=0.45\textwidth]{lecture-24/unbiased-var}}
    \end{center}
    The above is an image outlining the ideas of the proof
    for $n = 5$ and $k = 2$.
    \begin{itemize}
        \item The orange subspace is
        the column span of $\mathbb{X}$,
        with vector $\mathbb{X}\beta^*$ lying on it.
        \begin{itemize}
            \item We assume this subspace has dimension $k = 2$.
        \end{itemize}
        \item The vector $Y$ is the result of taking
        $\mathbb{X}\beta^*$ and adding $5$-dimensional Gaussian noise
        $\epsilon \sim \mathcal{N}_5(0, \sigma^2 I_5)$.
        \item The foot of the altitude of $Y$ onto the orange plane
        is $\mathbb{X}\hat{\beta}^{\mathrm{LS}}$,
        and $\hat{\epsilon}$ is the altitude vector itself.
        \item Crucially, only $3$ of the $5$ ``directions''
        of Gaussian noise point orthogonally to the orange
        $2$-dimensional subspace.
        \begin{itemize}
            \item It is only these $3$ directions of noise
            that contribute to the squared length of $\hat{\epsilon}$.
            \item Thus, we might expect
            the squared length of $\hat{\epsilon}$
            to be the sum of $3$ copies of $\sigma^2$.
        \end{itemize}
    \end{itemize}
    That's why we might expect
    $\mathbb{E}\left[\|\hat{\epsilon}\|^2\right] = 3\sigma^2$;
    in general, the coefficient of $3$
    would be $n - k$.
    The technical complexity of the proof
    comes from having to formalize the notion of
    ``directions of noise''.
\end{remark}

\begin{proof}
    The first step is to rewrite $\hat{\epsilon}$
    as the ``orthogonal component''
    of the overall Gaussian noise $\epsilon$.

    \begin{quote}
        \textbf{Lemma 1.}
        Let $H$ denote the $\mathbb{R}^n \to \mathbb{R}^n$ map
        given by projecting vectors onto the column space of $\mathbb{X}$.
        Then $\hat{\epsilon} = (I_n - H)\epsilon$.

        \emph{Proof:}
        Recalling from \href{/18.650/lectures/23/#interpreting-the-ls-solution}{Lecture 23}
        that $HY = \mathbb{X}\hat{\beta}^{\mathrm{LS}}$,
        we can rewrite the definition of $\hat{\epsilon}$ like so:
        \begin{align*}
            \hat{\epsilon}
        = Y - \mathbb{X}\hat{\beta}^{\mathrm{LS}}
        & = (I_n - H)Y \\
        & = (I_n - H)(\mathbb{X}\beta^* + \epsilon) \\
        & = (I_n - H)(\mathbb{X}\beta^*) +
        (I_n - H)\epsilon.
        \end{align*}
        However, $(I_n - H)(\mathbb{X}\beta^*) = 0$ because
        $\mathbb{X}\beta^*$ already lies on
        the column space of $\mathbb{X}$---that is,
        because $H(\mathbb{X}\beta^*) = \mathbb{X}\beta^*$.
        Then all that remains is the $\hat{\epsilon} = (I_n - H)\epsilon$
        term, as promised. ~ $\square$
    \end{quote}

    Define the matrix $M := I_n - H$;
    think of $M$ not as a projection map,
    but as an ``altitude'' map.
    Then we seek to show that
    $\mathbb{E}\left[\left\|M\epsilon\right\|^2\right] = (n - k) \sigma^2$,
    where $\epsilon \sim \mathcal{N}_n(0, \sigma^2 I_n)$.

    To proceed, we'll need some deeper linear algebra tricks; refer to
    \href{/18.701/lectures/18/#the-spectral-theorem}{18.701 Notes}
    for reference.

    \begin{quote}
        \textbf{Lemma 2.}
        The matrix $M$ diagonalizes as $M = U\Lambda U^{\top}$, where:
        \begin{itemize}
            \item The matrix $U$ is orthogonal.
            \item The matrix $\Lambda$ is diagonal
            with $k$ zeroes and $n - k$ ones.
        \end{itemize}

        \emph{Proof:}
        Note that $M$ is a real symmetric matrix,
        so the Spectral theorem applies.
        This, alone, already guarantees us
        that $M$ is expressible in the form
        $M = U \Lambda U^{\top}$
        with $U$ orthogonal and $\Lambda$ diagonal.

        It remains to study $\Lambda$, or equivalently,
        to study the spectrum of $M$.
        Recall that for the $k$-dimensional subspace
        $W := \mathrm{colspan}(\mathbb{X})$ of $\mathbb{R}^n$,
        the map $M$ takes vectors in $\mathbb{R}^n$
        to their ``altitudes'' onto $W$.  Therefore,
        \begin{itemize}
            \item Every vector along $W$
            has eigenvalue zero.
            \item Every vector along $W^{\perp}$
            (the $(n - k)$-dimensional subspace of vectors orthogonal to $W$)
            has eigenvalue one.
        \end{itemize}
        This implies the spectrum of $M$ contains
        $k$ zeroes and $(n - k)$ ones,
        proving the lemma's promise about $\Lambda$.
        ~ $\square$
    \end{quote}

    Denote $\eta := U^{\top}\epsilon$ (that is,
    $\epsilon$ expressed in the eigenbasis of $M$).
    Using Lemma 2, we can rewrite $\|M\epsilon\|^2$ like so:
    \[
    \| M\epsilon \|^2
    = \| U\Lambda U^{\top} \epsilon \|^2
    = \| \Lambda U^{\top} \epsilon \|^2
    = \| \Lambda \eta \|^2
    = \sum_{i = 1}^n \Lambda_{i, i}^2 \eta_i^2.
    \]
    Note that the step
    $\|U\Lambda U^{\top}\epsilon\|^2 = \|\Lambda U^{\top}\epsilon\|^2$
    comes from the orthogonality of $U$;
    rotating $\Lambda U^{\top}\epsilon$ via multiplication by $U$
    does not change its magnitude.
    Now just one more lemma.

    \begin{quote}
        \textbf{Lemma 3.}
        We have $\mathbb{E}[\eta] = 0$ and $\Cov(\eta) = \sigma^2 I_n$
        (the same moments as $\epsilon$).
        By corollary, $\mathbb{E}[\eta_i^2] = \sigma^2$ for all $i$.

        \emph{Proof:}
        The point is that $U$ is orthogonal,
        so $U^{\top}U = I_n$, making the covariance unchanged.
        \[ \mathbb{E}[\eta] = U^{\top} \mathbb{E}[\epsilon] = 0
        ~~~ \text{ and } ~~~
        \Cov(\eta)
        = U^{\top}\Cov(\epsilon)U 
        = \sigma^2 I_n \cdot U^{\top}U
        = \sigma^2 I_n. ~~~ \square \]
    \end{quote}
    By Lemma 3, taking expectations on both sides yields:
    \[ \mathbb{E}\left[
        \|M\epsilon\|^2
    \right]
    = \mathbb{E}\left[
        \sum_{i = 1}^n \Lambda_{i, i}^2 \eta_i^2
    \right]
    = \sum_{i = 1}^n \Lambda_{i, i}^2 \mathbb{E}\left[\eta_i^2\right]
    = \sum_{i = 1}^n \Lambda_{i, i}^2 \sigma^2. \]
    Exactly $n - k$ terms have $\Lambda_{i, i} = 1$,
    and the rest have $\Lambda_{i, i} = 0$,
    so the summation evaluates to $(n - k) \sigma^2$,
    as promised. ~ $\blacksquare$
\end{proof}

With all that out of the way\dots
we can now construct confidence intervals
and perform hypothesis tests
with $\hat{\beta}^{\mathrm{LS}}$.

\begin{remark}
    Most of the time,
    these tests will ask for whether
    $\beta_i^* = 0$ or $\beta_i^* \neq 0$
    for some particular entry $i$
    of $\beta^*$.
    This translates to asking
    whether the information communicated
    by the $i^{\text{th}}$ entry of a feature vector
    says anything about the response
    that the other entries do not already say.
\end{remark}

\subsection{Affine Linear Regression}

Suppose I have a collection of data points
$\{(X_i, Y_i)\}_{i = 1}^n$
on the coordinate plane,
and I'd like to draw their line of best fit.
In other words, I would like to find
parameters $(\hat{\beta}_1^{\mathrm{LS}}, \hat{\beta}_0^{\mathrm{LS}})$
for which $Y_i \approx \hat{\beta}^{\mathrm{LS}}_1X_i
+ \hat{\beta}^{\mathrm{LS}}_0$.

This is \emph{not} the same as
a one-dimensional linear regression,
as that would only produce a model
$Y_i \approx \hat{\beta}^{\mathrm{LS}}X_i$.
Fortunately, there's an easy adjustment
we can make to linear regression
to generalize it to affine models as well.

\textbf{Idea.}
Pad every feature vector $X_i \in \mathbb{R}^k$
with an extra $1$, extending it to $\mathbb{R}^{k + 1}$.  Then:
\[
\underbrace{
\begin{pmatrix}
Y_1 \\ Y_2 \\ \vdots \\ Y_n
\end{pmatrix}
}_{Y}
\approx
\underbrace{
\begin{pmatrix}
1 & \text{---} & X_1^{\top} & \text{---} \\
1 & \text{---} & X_2^{\top} & \text{---} \\
\vdots & & \vdots & \\
1 & \text{---} & X_n^{\top} & \text{---} \\
\end{pmatrix}
}_{\mathbb{X}}
\underbrace{
\begin{pmatrix}
\beta_0 \\ \beta_1 \\ \vdots \\ \beta_k
\end{pmatrix}
}_{\hat{\beta}^{\mathrm{LS}}}
=
\begin{pmatrix}
\text{---} & X_1^{\top} & \text{---} \\
& \vdots & \\
\text{---} & X_n^{\top} & \text{---} \\
\end{pmatrix}
\begin{pmatrix}
\beta_1 \\ \vdots \\ \beta_k
\end{pmatrix}
+ \begin{pmatrix}
\beta_0 \\ \vdots \\ \beta_0
\end{pmatrix}.
\]
So drawing a line of best fit
when the feature dimension is $k = 1$
is secretly a two-dimensional linear regression,
where the first entry of $\hat{\beta}^{\mathrm{LS}}$
communicates the $y$-intercept,
and the second entry communicates the slope.

\subsection{1D Linear Regression and Correlation}

\textbf{Definition.}
Given \emph{one-dimensional} data
$X_1, \dots, X_n \in \mathbb{R}$
and $Y_1, \dots, Y_n \in \mathbb{R}$,
we may compute their
\textbf{sample variances} and \textbf{sample covariances}
like so:
\[ S_{xx} := \frac{1}{n} \sum_{i = 1}^n (X_i - \bar{X}_n)^2
~~~ \parallel ~~~
S_{yy} := \frac{1}{n} \sum_{i = 1}^n (Y_i - \bar{Y}_n)^2
~~~ \parallel ~~~
S_{xy} := \frac{1}{n} \sum_{i = 1}^n (X_i - \bar{X}_n)(Y_i - \bar{Y}_n) \]

\textbf{Definition.}
The \textbf{correlation coefficient}
between the $\{X_i\}_{i = 1}^n$ and the $\{Y_i\}_{i = 1}^n$
is defined as $r := \frac{S_{xy}}{\sqrt{S_{xx}S_{yy}}}$.

It turns out that in the case of
one-dimensional line-of-best-fit (or 2D linear regression),
we have the following:

\begin{theorem}{1D Linear Regression Slope}
    In a one-dimensional line-of-best-fit
    $Y_i \approx \hat{\beta}^{\mathrm{LS}}_1X_i + \hat{\beta}^{\mathrm{LS}}_0$,
    we have $\hat{\beta}^{\mathrm{LS}}_1 = \frac{S_{xy}}{S_{xx}}$.
\end{theorem}

\begin{proof}
    Details omitted; the proof is just unenlightening computation.
    ~ $\blacksquare$
\end{proof}

Now let $\hat{Y}_i :=
\hat{\beta}^{\mathrm{LS}}_1X_i + \hat{\beta}^{\mathrm{LS}}_0$
denote the value of $Y_i$ predicted by the line-of-best-fit
(in contrast to the true value $Y_i$).

\begin{theorem}{Correlation and Regression}
    We have that
    $\frac{\hat{Y}_i - \bar{Y}_n}{\sqrt{S_{yy}}}
    = r \cdot \frac{X_i - \bar{X}_n}{\sqrt{S_{xx}}}$.
\end{theorem}

\begin{proof}
    More computation; omitted. ~ $\blacksquare$
\end{proof}

And that's why AP Statistics makes you
interpret the correlation coefficient by writing,
``The constant $r^2$ represents
the percent of the variation in the response
that is explained by the linear relationship with
the feature vector.''

\end{document}